\documentclass[12pt]{article}
\usepackage{tcolorbox}

\include{preamble}

\newtoggle{professormode}
%\toggletrue{professormode} %STUDENTS: DELETE or COMMENT this line



\title{MATH 342W / 642 / RM742 Spring \the\year~HW \#1}

\author{Bryan Nevarez} %STUDENTS: write your name here

\iftoggle{professormode}{
\date{Due 11:59PM February 11, \the\year~by email \\ \vspace{0.5cm} \small (this document last updated \currenttime~on \today)}
}

\renewcommand{\abstractname}{Instructions and Philosophy}

\begin{document}
\maketitle

\iftoggle{professormode}{
\begin{abstract}
The path to success in this class is to do many problems. Unlike other courses, exclusively doing reading(s) will not help. Coming to lecture is akin to watching workout videos; thinking about and solving problems on your own is the actual ``working out.''  Feel free to \qu{work out} with others; \textbf{I want you to work on this in groups.}

Reading is still \textit{required}. For this homework set, read the first chapter of \qu{Learning from Data} and the introduction and Chapter 1 of Silver's book. Of course, you should be googling and reading about all the concepts introduced in class online. This is your responsibility to supplement in-class with your own readings.

The problems below are color coded: \ingreen{green} problems are considered \textit{easy} and marked \qu{[easy]}; \inorange{yellow} problems are considered \textit{intermediate} and marked \qu{[harder]}, \inred{red} problems are considered \textit{difficult} and marked \qu{[difficult]} and \inpurple{purple} problems are extra credit. The \textit{easy} problems are intended to be ``giveaways'' if you went to class. Do as much as you can of the others; I expect you to at least attempt the \textit{difficult} problems. 

This homework is worth 100 points but the point distribution will not be determined until after the due date. See syllabus for the policy on late homework.

Up to 7 points are given as a bonus if the homework is typed using \LaTeX. Links to instaling \LaTeX~and program for compiling \LaTeX~is found on the syllabus. You are encouraged to use \url{overleaf.com}. If you are handing in homework this way, read the comments in the code; there are two lines to comment out and you should replace my name with yours and write your section. The easiest way to use overleaf is to copy the raw text from hwxx.tex and preamble.tex into two new overleaf tex files with the same name. If you are asked to make drawings, you can take a picture of your handwritten drawing and insert them as figures or leave space using the \qu{$\backslash$vspace} command and draw them in after printing or attach them stapled.

The document is available with spaces for you to write your answers. If not using \LaTeX, print this document and write in your answers. I do not accept homeworks which are \textit{not} on this printout. Keep this first page printed for your records.

\end{abstract}

\thispagestyle{empty}
\vspace{1cm}
NAME: \line(1,0){380}
\clearpage
}

%============================================
% PROBLEM #1: Nate Silver's The Signal and Noise (Introduction & Chapter 1)
%============================================
\problem{These are questions about Silver's book, the introduction and chapter 1.}

\begin{enumerate}

%--------------------------------------------
% Part (a)
%--------------------------------------------
\easysubproblem{What is the difference between \emph{predict} and \emph{forecast}? Are these two terms used interchangably today?}\spc{4}

\begin{tcolorbox}[colback=white, sharp corners]

These two terms are used interchangeably today. The term \emph{predict} is the act of making a claim/statement regarding a future event. In slight contrast, the term \emph{forecast}, according to Silver, was used to plan for the future, using ``prudence, wisdom, and industriousness,'' under uncertain conditions. The term ``forecast'' shared more similarity with how we use the term ``foresight'' than with \emph{predict}.


\end{tcolorbox}

%--------------------------------------------
% Part (b)
%--------------------------------------------
\easysubproblem{What is John P. Ioannidis's findings and what are its implications?}\spc{5}

\begin{tcolorbox}[colback=white, sharp corners]

Silver mentions the findings of John P. Iaonnidis from his controversial paper entitled \emph{``Why Most Pubhlished Reserach Findings are False''} as a series of ``pitfalls'' where big data has not measured up with its proposed utility. According to Iaonnidis's findings, most of the successful predictions of medical hypotheses carried out in labs written in peer-reviewed journals were in fact not reproducible as verified by a third-party, Bayer Laboratories. The implication is that big data is not without its shortcomings and may not be fully reliable nor reproducible of its effectiveness. 

\end{tcolorbox}

%--------------------------------------------
% Part (c)
%--------------------------------------------
\easysubproblem{What are the human being's most powerful defense (according to Silver)? Answer using the language from class.}\spc{4}

\begin{tcolorbox}[colback=white, sharp corners]

Human beings' most superior defensive attribute, compared to the rest of the living organism that we share the planet with, is the ability to ``generalize'' which would be the ability to find/identify patterns among phenomena. In other words, we as a species have an innate tendency to find what $f$ is and in certain contexts we are good at it. 

\end{tcolorbox}

%--------------------------------------------
% Part (d)
%--------------------------------------------
\easysubproblem{Information is increasing at a rapid pace, but what is not increasing?}\spc{3}

\begin{tcolorbox}[colback=white, sharp corners]

The quantity and rate at which information grows each day has been increasing (and will continue to do so), however, ``\emph{useful} information'' has not increased at the same rate. So most of this growth of information is categorized as ``\emph{noise}.''

\end{tcolorbox}

%--------------------------------------------
% Part (e)
%--------------------------------------------
\hardsubproblem{Silver admits that we will always be subjectively biased when making predictions. However, he believes there is an objective truth. In class, how did we describe the objective truth? Answer using notation from class i.e. $t,f, g, h^*, \delta, \epsilon, t, z_1, \ldots, z_t, \delta, \mathbb{D}$, $\mathcal{H}, \mathcal{A}, \mathcal{X}, \mathcal{Y}, X, y, n, p$, $x_{\cdot 1}, \ldots, x_{\cdot p}, x_{1 \cdot}, \ldots, x_{n \cdot}$, etc.}\spc{3}

\begin{tcolorbox}[colback=white, sharp corners]

The \emph{objective truth} of a phenomenon is encapsulated by the following terms and notation:

\begin{itemize}
\item proximal causes which are the true drivers of the phenomenon, denoted by $z_1, \ldots, z_t$
\item the unknown function that takes the $z_i$'s as its inputs, denoted by $t$
\item the response/output/dependent variable/regressand/outcome/endpoint denoted by $y$
\end{itemize}

\end{tcolorbox}

%--------------------------------------------
% Part (f)
%--------------------------------------------
\easysubproblem{In a nutshell, what is Karl Popper's (a famous philosopher of science) definition of \emph{science}?}\spc{4}

\begin{tcolorbox}[colback=white, sharp corners]

According to Karl Popper, the definition of \emph{science} is the continual update/testing of a theory/theories. The tests determine which theory needs to be refuted and updated based upon how well the predicted result held up in the real world. 

\end{tcolorbox}

%--------------------------------------------
% Part (g)
%--------------------------------------------
\intermediatesubproblem{Why did the ratings agencies say the probability of a CDO defaulting was 0.12\% instead of the 28\% that actually occured? Answer using concepts from class.}\spc{4}

\begin{tcolorbox}[colback=white, sharp corners]

Silver delves into the crux of the agencies faulty ratings as to whether it was due to avarice or ignorance. The agencies ratings was based upon a predictive model that made the incorrect assumption that the risks were ``well-diversified.'' Using the concepts/terminology discussed in class, the agencies' model assumed that the proxies to the proximal causes, denoted by $x_1, \ldots, x_n$, were all independent from each other. As a result, the predicted risk for the mortgages to default were off by a factor of 160,000. 

\end{tcolorbox}

%--------------------------------------------
% Part (h)
%--------------------------------------------
\easysubproblem{What is the difference between \emph{risk} and \emph{uncertainty} according to Silver's definitions?}\spc{4}

\begin{tcolorbox}[colback=white, sharp corners]

According to Silver's definitions, \emph{risk} is an exact calculated likelihood, e.g. the odds, of an occurence of an event (usually often characterized as an unwanted, deleterious outcome). On the other hand, \emph{uncertainty} is the same thing as risk just without the numerical value, i.e., it is risk that cannot be measured/calculated/tabulated. 

\end{tcolorbox}

%--------------------------------------------
% Part (i)
%--------------------------------------------
\hardsubproblem{How does Silver define \emph{out of sample}? Answer using notation from class i.e. $t,f, g, h^*, \delta, \epsilon, z_1, \ldots, z_t, \delta, \mathbb{D}, \mathcal{H}, \mathcal{A}, \mathcal{X}, \mathcal{Y}, X, y, n, p, x_{\cdot 1}, \ldots, x_{\cdot p}, x_{1 \cdot}, \ldots, x_{n \cdot}$, etc. WARNING: Silver defines \emph{out of sample} completely differently than the literature, than practitioners in industry and how we will define it in class in a month or so. We will explore what he is talking about in class in the future and we will term this concept differently, using the more widely accepted terminology. So please forget the phrase \emph{out of sample} for now as we will introduce it later in class as something else. There will be other such terms in his book and I will provide this disclaimer at these appropriate times.}\spc{7}

\begin{tcolorbox}[colback=white, sharp corners]

Silver defines an \emph{out-of-sample} problem as a situation where one does not have data any historical data on a potential risky activity or result. Therefore, the model does not have any way to forecast said risk and the predictions do not account for the potential deleterious results of the risky behavior. In other words, there exists elements in our input space $\mathcal{X}$ that are not part of the training data $\mathbb{D}$. There would need to be additional features $x_{\cdot p+1}, x_{\cdot p+1}, \ldots $ that would account for the situations that would be considered \emph{out-of-sample} if they were not included in the measurement/data collection process.. Without these additional features for the data collection/measurement of the real phenomenon our $g$ would not be able to even forecast/predict anything related to those \emph{out-of-sample} events. 
\end{tcolorbox}

%--------------------------------------------
% Part (j)
%--------------------------------------------
\intermediatesubproblem{Look up \emph{bias} and \emph{variance} online or in a statistics textbook. Connect these concepts to Silver's terms \emph{accuracy} and \emph{precision}. This is another example of Silver using non-standard terminology.}\spc{9}

\begin{tcolorbox}[colback=white, sharp corners]
\emph{Bias} is error based upon the erroneous nature of the assumptions of the model which may result in \emph{underfitting} the data. \emph{Variance} is error based upon the sensitivity to the fluctuations of the training data, a.k.a. noise.
\emph{Bias} is akin to Silver's notion of \emph{accuracy} and \emph{variance} is analagous to Silver's notion of \emph{precision}. The lower the bias of a model the more accurate it is, the lower the variance and the better the precision of the model. 

\end{tcolorbox}

\end{enumerate}

\newpage
%============================================
% PROBLEM #2: Theory of Modeling
%============================================
\problem{Below are some questions about the theory of modeling.}

\begin{enumerate}

%--------------------------------------------
% Part (a)
%--------------------------------------------
\easysubproblem{Redraw the illustration of Earth and the table-top globe except do not use the Earth and a table-top globe as examples (use another example). The quadrants are connected with arrows. Label these arrows appropriately.}\spc{11}

\begin{tcolorbox}[colback=white, sharp corners]

\vspace{-1.8in}

\begin{center}
\includegraphics[scale=.7]{modeling}
\end{center}

\vspace{-1.8in}

\end{tcolorbox}

%--------------------------------------------
% Part (b)
%--------------------------------------------
\easysubproblem{Pursuant to the fix in the previous question, how do we define \emph{data} for the purposes of this class?}\spc{2}

\begin{tcolorbox}[colback=white, sharp corners]

For the purposes of this class, \emph{data} is the natural result of measuring the phenomenon (weather). 

\end{tcolorbox}

%--------------------------------------------
% Part (c)
%--------------------------------------------
\easysubproblem{Pursuant to the fix in the previous question, how do we define \emph{predictions} for the purposes of this class?}\spc{3}

\begin{tcolorbox}[colback=white, sharp corners]

For the purposes of this class, \emph{predictions} are what the model tells us what will happen in a certain future (meteorological report). 

\end{tcolorbox}

%--------------------------------------------
% Part (d)
%--------------------------------------------
\easysubproblem{Why are \qu{all models wrong}? We are quoting the famous statisticians George Box and Norman Draper here.}\spc{2}

\begin{tcolorbox}[colback=white, sharp corners]

``All models'' are considered to be ``wrong'' because they are not precise nor exactly correct with predicting the behavior of a particular phenomena of interest. By definition, models are approximations. Much of the phenomena in the physical world we live in are perturbations from the exactitude of an analytical solution to a theoretical problem that describes that phenomena. 

\end{tcolorbox}

%--------------------------------------------
% Part (e)
%--------------------------------------------
\intermediatesubproblem{Why are \qu{[some models] useful}? We are quoting the famous statisticians George Box and Norman Draper here.}\spc{2}

\begin{tcolorbox}[colback=white, sharp corners]

Despite all models ``being wrong,'' some are considered to be useful because some of the models provide a reasonable degree of accuracy and consistency with its description of the phenomena of interest. The findings from these useful models may be utilized to further comprehend the phenomena and to even predict future behavior. 

\end{tcolorbox}

%--------------------------------------------
% Part (f)
%--------------------------------------------
\intermediatesubproblem{What is the difference between a "good model" and a "bad model"?}\spc{2}

\begin{tcolorbox}[colback=white, sharp corners]

A ``bad model'' does makes inaccurate and/or imprecise predictions upon new data. A ``good model'' minimizes the error made between the predicted value and the true value of new data. A ``good'' model also strikes an optimal balance between its bias (``underfitting'') and variance (``overfitting'') of the training data.

\end{tcolorbox}

\end{enumerate}


\newpage
%============================================
% PROBLEM #3: Modeling Aphorisms
%============================================
\problem{We are now going to investigate the famous English aphorism \qu{an apple a day keeps the doctor away} as a model. We will use this as springboard to ask more questions about the framework of modeling we introduced in this class.}

\begin{enumerate}

%--------------------------------------------
% Part (a)
%--------------------------------------------
\easysubproblem{Is this a mathematical model? Yes / no and why.}\spc{3}

\begin{tcolorbox}[colback=white, sharp corners]

The English aphorism above can be represented by a mathematical model because there is an inherent counting aspect in the aphorism which can be measured and therefore metrics can be constructed and implemented to model and predict the efficaciousness of the proposed solution for ``keeping the doctor away.''

\end{tcolorbox}

%--------------------------------------------
% Part (b)
%--------------------------------------------
\easysubproblem{What is(are) the input(s) in this model?}\spc{3}

\begin{tcolorbox}[colback=white, sharp corners]

The input would be the number of days where one would eat an apple. 

\end{tcolorbox}

%--------------------------------------------
% Part (c)
%--------------------------------------------
\easysubproblem{What is(are) the output(s) in this model?}\spc{3}

\begin{tcolorbox}[colback=white, sharp corners]

The output would be how long would it be until you have to go to the doctor under the assumption that you go to the doctor because of an ailment/health issue.

\end{tcolorbox}

%--------------------------------------------
% Part (d)
%--------------------------------------------
\intermediatesubproblem{How good / bad do you think this model is and why?}\spc{3}

\begin{tcolorbox}[colback=white, sharp corners]

This may not be a great model because there presumably are many more reasons that would affect whether or not one seeks medical attention from a doctor other than simply knowing whether or not one eats an apple a day. 

\end{tcolorbox}

%--------------------------------------------
% Part (e)
%--------------------------------------------
\easysubproblem{Devise a metric for gauging the main input. Call this $x_1$ going forward.}\spc{4}

\begin{tcolorbox}[colback=white, sharp corners]

A metric I would propose to gauge the main input would be 

\vspace{-5mm}

\[
x_1 = \text{the avg. number of days that an apple was eaten per year b/w ages 5 - 60}
\]

\end{tcolorbox}

%--------------------------------------------
% Part (f)
%--------------------------------------------
\easysubproblem{Devise a metric for gauging the main output. Call this $y$ going forward.}\spc{4}

\begin{tcolorbox}[colback=white, sharp corners]

For gauging the main output, I would propose to measure 

\vspace{-5mm}

\[
y = \text{the number of days (time) between the last day that the person eats}\\
\]
\[
\text{ an apple and the first time the person goes to see a doctor for medical treatment}
\]

\end{tcolorbox}

%--------------------------------------------
% Part (g)
%--------------------------------------------
\easysubproblem{What is $\mathcal{Y}$ mathematically?}\spc{3}

\begin{tcolorbox}[colback=white, sharp corners]

\[
	\mathcal{Y} = \lbrace y \in \naturals_0 \mid  0 \leq y \leq 36500 \rbrace
\]

\end{tcolorbox}

%--------------------------------------------
% Part (h)
%--------------------------------------------
\easysubproblem{Briefly describe $z_1, \ldots, z_t$ in English where $y = t(z_1, \ldots, z_t)$ in this \emph{phenomenon} (not \emph{model}).}\spc{3}

\begin{tcolorbox}[colback=white, sharp corners]

The $z_1, \ldots, z_t$ are the proximal causes which are the true drivers for the phenomenon of how eating apples on a daily basis prevents a visit to the doctor. In reality, we do not know what these causes are. 

\end{tcolorbox}

%--------------------------------------------
% Part (i)
%--------------------------------------------
\easysubproblem{From this point on, you only observe $x_1$. What is the value of $p$?}\spc{1}

\begin{tcolorbox}[colback=white, sharp corners]

For this case, the value of $p$ is 1 because we only have one metric for one only proxy.

\end{tcolorbox}

%--------------------------------------------
% Part (j)
%--------------------------------------------
\intermediatesubproblem{What is $\mathcal{X}$ mathematically? If your information contained in $x_1$ is non-numeric, you must coerce it to be numeric at this point.}\spc{3}

\begin{tcolorbox}[colback=white, sharp corners]

\[
	\mathcal{X} = \lbrace x \in \rationals \mid 0\leq x \leq 365 \rbrace 
\]

\end{tcolorbox}

%--------------------------------------------
% Part (k)
%--------------------------------------------
\easysubproblem{How did we term the functional relationship between $y$ and $x_1$? Is it approximate or equals?}\spc{3}

\begin{tcolorbox}[colback=white, sharp corners]

\[
	y \neq f(x_1) \implies y \approx f(x_1)
\]

\end{tcolorbox}

%--------------------------------------------
% Part (l)
%--------------------------------------------
\easysubproblem{Briefly describe \emph{superivised learning}.}\spc{5}

\begin{tcolorbox}[colback=white, sharp corners]

\emph{Supervised learning} is the process modeling a particularly phenomenon of interest by finding the best function $g: \mathcal{X} \to \mathcal{Y}$ that approximates the ideal, unknown target function $f: \mathcal{X} \to \mathcal{Y}$ where $\mathcal{X}$ is the input space and $\mathcal{Y}$ is the output space through the application of an algorithm that arrives upon its solution through iterating over training data $\mathbb{D}$ that has known correct output values $y \in \mathcal{Y}$ for given $\mathbf{x} \in \mathcal{X}$.

\end{tcolorbox}

\newpage

%--------------------------------------------
% Part (m)
%--------------------------------------------
\easysubproblem{Why is \emph{superivised learning} an \emph{empirical solution} and not an \emph{analytic solution}?}\spc{3}

\begin{tcolorbox}[colback=white, sharp corners]

\emph{Supervised learning} is an approach of ``learning from [historical] data'' and henceforth it produces an \emph{empirical solution} and not an \emph{analytic solution.}

\end{tcolorbox}

%--------------------------------------------
% Part (n)
%--------------------------------------------
\intermediatesubproblem{From this point on, assume we are involved in supervised learning to achieve the goal you stated in the previous question. Briefly describe what $\mathbb{D}$ would look like here.}\spc{3}

\begin{tcolorbox}[colback=white, sharp corners]

$\mathcal{D}$ represents the training data, $X$ with the known output values in $\mathbf{y}$.

\end{tcolorbox}

%--------------------------------------------
% Part (o)
%--------------------------------------------
\intermediatesubproblem{Briefly describe the role of $\mathcal{H}$ and $\mathcal{A}$ here.}\spc{3}

\begin{tcolorbox}[colback=white, sharp corners]

The set of candidate functions, denoted by $\mathcal{H}$, contains elements $h$ that approximate the function $f$. This set is useful because it restricts the space of all functions into a more well-defined set and helps with finding the ``best one'' that would model the phenomenon.
$\mathcal{A}$ is the algorithm that takes the training set data, $\mathbb{D}$ and $\mathcal{H}$ and returns $g$ which is an approximation to $f$, i.e. $g = \mathcal{A}\left(\mathbb{D}, \mathcal{H} \right)$.


\end{tcolorbox}

%--------------------------------------------
% Part (p)
%--------------------------------------------
\easysubproblem{If $g = \mathcal{A}(\mathbb{D}, \mathcal{H})$, what should the domain and range of $g$ be?}\spc{3}

\begin{tcolorbox}[colback=white, sharp corners]

Domain of $g$: $\lbrace x \in \rationals \mid 0\leq x \leq 365 \rbrace $ \quad Range of $g$: $\lbrace y \in \naturals_0 \mid  0 \leq y \leq 36500 \rbrace$

\end{tcolorbox}

%--------------------------------------------
% Part (q)
%--------------------------------------------
\easysubproblem{Is $g \in \mathcal{H}$? Why or why not?}\spc{3}

\begin{tcolorbox}[colback=white, sharp corners]

Yes, $g \in \mathcal{H}$ because $g$ comes from the algorithm $\mathcal{A}$ after going through the functions in $\mathcal{H}$ and seeing which one produces the minimal error based upon the training data $\mathcal{D}$.
\end{tcolorbox}

%--------------------------------------------
% Part (r)
%--------------------------------------------
\easysubproblem{Given a never-before-seen value of $x_1$ which we denote $x^*$, what formula would we use to predict the corresponding value of the output? Denote this prediction $\hat{y}^*$.}\spc{1}

\begin{tcolorbox}[colback=white, sharp corners]

For a new, never-before-seen value of $x_1$ which we denote $x^{*}$, we will represent the corresponding value of the output by the following, 

\[
	\hat{y} = g(x^{*})
\]

\end{tcolorbox}

%--------------------------------------------
% Part (s)
%--------------------------------------------
\intermediatesubproblem{$f$ is the function that is the best possible fit of the phenomenon given the covariates. We will unfortunately not be able to define \qu{best} until later in the course. But you can think of it as a device that extracts all possible information from the covariates and whatever is left over $\delta$ is due exclusively to information you do not have. Is it reasonable to assume $f \in \mathcal{H}$? Why or why not?}\spc{2}

\begin{tcolorbox}[colback=white, sharp corners]

No, it is not reasonable to assume that $f \in \mathcal{H}$ because $f$ is assumed \emph{a priori} that $f$ is arbitrarily complicated and unknown while $\mathcal{H}$ usually contains much more simplistic and known functions which can be fit by the application of a particular algorithm, $\mathcal{A}$. 

\end{tcolorbox}

%--------------------------------------------
% Part (t)
%--------------------------------------------
\easysubproblem{In the general modeling setup, if $f \notin \mathcal{H}$, what are the three sources of error? Copy the equation from the class notes. Denote the names of each error and provide a sentence explanation of each. Denote also $e$ and $\mathcal{E}$ using underbraces / overbraces.}\spc{6}

\begin{tcolorbox}[colback=white, sharp corners]

We have three types of errors:

\begin{itemize}
	\item[1.] estimation error = $h^{*}(x_1, \ldots, x_p) - g(x_1, \ldots, x_p)$\\
	Estimation error determines how good or bad the algorithm is in choosing the modeling function $g$ compared with the candidate model, $h^{*}$ from $\mathcal{H}$.
	\item[2.] misspecification error = $f(x_1, \ldots, x_p) - h^{*}(x_1, \ldots, x_p)$\\
	Misspecification error measures is the inherent error that arises from the assumptions that are made in making the modeling claims to the phenomenon of interest.
	\item[3.] error due to ignorance: $\delta = t - f$\\
	This error is a result of the built-in ignorance we have of the true causal drivers for the phenomena we are modeling. The proxies $x_1, \ldots, x_p$ are not the true causal drivers $z_1, \ldots, z_p$ and so we are missing information. 
\end{itemize}

Let $\mathbf{x} = (x_1, \ldots, x_p)$. We then have the following, 

\[
	y  = g(\mathbf{x}) + \underbrace{\left(h^{*}(\mathbf{x}) - g(\mathbf{x}) \right) - \overbrace{\left(f(\mathbf{x}) - h^{*}(\mathbf{x}) \right) + \delta}^{\mathcal{E}}}_{e}
\]

\end{tcolorbox}

\newpage
%--------------------------------------------
% Part (u)
%--------------------------------------------
\easysubproblem{In the general modeling setup, for each of the three source of error, explain what you would do to reduce the source of error as best as you can.}\spc{7}

\begin{tcolorbox}[colback=white, sharp corners]

\begin{enumerate}
	\item[1.] To decrease estimation error we need to increase the sample size, $n$, from more historical examples which are the rows in $\mathbb{D}$.
	\item[2.] We decrease the model misspecification error by expanding the set of candidate functions in $\mathcal{H}$ to be more complex with the goal of encapsulating the complex relationships among the proxies better. 
	\item[3.] We decrease the error due to ignorance by increasing $p$ with more useful variables. 
\end{enumerate}

\end{tcolorbox}

%--------------------------------------------
% Part (v)
%--------------------------------------------
\intermediatesubproblem{In the general modeling setup, make up an $f$, an $h^*$ and a $g$ and plot them on a graph of $y$ vs $x$ (assume $p=1$). Indicate the sources of error on this plot (see last question). Which source of error is missing from the picture? Why?}\spc{9}

\begin{tcolorbox}[colback=white, sharp corners]

The source of error that is missing from the picture is $\delta$, the error due to ignorance, because we do not have any information on the ``true'' function $t$. 

\end{tcolorbox}

%--------------------------------------------
% Part (w)
%--------------------------------------------
\easysubproblem{What is a null model $g_0$? What data does it make use of? What data does it not make use of?}\spc{3}

\begin{tcolorbox}[colback=white, sharp corners]

The ``null model'' $g_0$ is the model we end up with if we do not have any $x_i$'s to ``learn'' from and we only use the data from the known $y$'s.  

\end{tcolorbox}

%--------------------------------------------
% Part (x)
%--------------------------------------------
\easysubproblem{What is a parameter in $\mathcal{H}$?}\spc{3}

\begin{tcolorbox}[colback=white, sharp corners]

A parameter in $\mathcal{H}$ represents a coefficient that varies among the set of candidate function in $\mathcal{H}$.

\end{tcolorbox}

%--------------------------------------------
% Part (y)
%--------------------------------------------
\easysubproblem{Regardless of your answer to what $\mathcal{Y}$ was above in (g), we now coerce $\mathcal{Y} = \braces{0,1}$. What would the null model $g_0$ be and why?}\spc{2}

\begin{tcolorbox}[colback=white, sharp corners]

\[
	g_0 = \text{Mode}\left[ \mathbf{y} \right] \in \lbrace 0, 1 \rbrace
\]

This is as simplistic of a model as we can get without any training data.

\end{tcolorbox}

\newpage
%--------------------------------------------
% Part (z)
%--------------------------------------------
\easysubproblem{Regardless of your answer to what $\mathcal{Y}$ was above in (g), we now coerce $\mathcal{Y} = \braces{0,1}$. If we use a threshold model, what would $\mathcal{H}$ be? What would the parameter(s) be?}\spc{2}

\begin{tcolorbox}[colback=white, sharp corners]

Given the $\mathcal{Y} = \braces{0,1}$ and if we are using a threshold model, we will have the following: 

\[
	\mathcal{H} = \lbrace \indic{x \geq \theta} : \theta \in \mathcal{X} \rbrace
\]

In this scenario, $\theta$ is the model parameter.

\end{tcolorbox}

%--------------------------------------------
% Part (aa)
%--------------------------------------------
\easysubproblem{Give an explicit example of $g$ under the threshold model.}\spc{1}

\begin{tcolorbox}[colback=white, sharp corners]

Let's suppose that threshold parameter $\theta = 200$. This means that for those people who eat an apple on, at least, an average of 200 out of the 365 days in a year over the course of ages 5 - 60 years will never have to go to visit the doctor. In this case, after we coerced $\mathcal{Y} = \lbrace 0, 1 \rbrace$, the output value 0 means that the person had to visit the doctor and 1 means that the person never had to visit the doctor. 

So, under this specific threshold model, if $x^{*}= 250$, then $g(x^{*}) = 1$. 

\end{tcolorbox}

\end{enumerate}

\newpage
%============================================
% PROBLEM #4: Richard Feynman lecture 
%============================================
\problem{As alluded to in class, modeling is synonymous with the entire enterprise of science. 

In 1964, \href{https://en.wikipedia.org/wiki/Richard_Feynman}{Richard Feynman}, a famous physicist and public intellectual with an inimitably captivating presentation style, gave a series of seven lectures in 1964 at Cornell University on the \qu{character of physical law}. Here is a \href{https://www.youtube.com/watch?v=EYPapE-3FRw}{10min excerpt} of one of these lectures about the scientific method. Feel free to watch the entire clip, but for the purposes of this class, we are only interested in the following segments: 0:00-1:00 and 3:48-6:45. }

\begin{enumerate}

%--------------------------------------------
% Part (a)
%--------------------------------------------
\intermediatesubproblem{According to Feynman, how does the scientific method differ from learning from data with regards to building models for reality? (0:08)}\spc{3}

\begin{tcolorbox}[colback=white, sharp corners]

The scientific method differs from learning from data with regards to building models for reality in the sense that in the scientific method model is ``guessed.''

\end{tcolorbox}

%--------------------------------------------
% Part (b)
%--------------------------------------------
\intermediatesubproblem{He uses the phrase \qu{compute consequences}. What word did we use in class for \qu{compute consequences}? This word also appears in your diagram in 2a. (0:14)}\spc{1}

\begin{tcolorbox}[colback=white, sharp corners]

We use the word ``predictions'' for ``compute consequences.''

\end{tcolorbox}

%--------------------------------------------
% Part (c)
%--------------------------------------------
\intermediatesubproblem{When he says compare consequences to \qu{experiment}, what word did we use in class for \qu{experiment}? This word also appears in your diagram in 2a. (0:29)}\spc{1}

\begin{tcolorbox}[colback=white, sharp corners]

For the word ``experiment'' that Feynman uses, we use the word ``measurement'' or ``data collection'' to mean the same thing. 

\end{tcolorbox}

%--------------------------------------------
% Part (d)
%--------------------------------------------
\intermediatesubproblem{When he says \qu{compare consequences to experiment}, which part of the diagram in 2a is that comparison?}\spc{1}

\begin{tcolorbox}[colback=white, sharp corners]

The 	``compare consequences to experiment'' is part of the diagram in 2a where we perform 	``model validation'' where we compare the predicted values to the measurements/data collected.

\end{tcolorbox}

%--------------------------------------------
% Part (e)
%--------------------------------------------
\hardsubproblem{When he says \qu{if it disagrees with experiment, it's wrong} (0:44), would a data scientist agree/disagree? What would the data scientist further comment?}\spc{3}

\begin{tcolorbox}[colback=white, sharp corners]

A data scientist may agree with the claim that the guess is ``wrong'' because it ``disagrees with experiment,'' however, the data scientist would further add that ``it may be useful.''

\end{tcolorbox}

%--------------------------------------------
% Part (f)
%-------------------------------------------
\hardsubproblem{[You can skip his UFO discussion as it belongs in a class on statistical inference on the topic of $H_0$ vs $H_a$ which is \emph{not} in the curriculum of this class.] He then goes on to say \qu{We can disprove any definite theory. We never prove [a theory] right... We can only be sure we're wrong} (3:48 - 5:08). What does this mean about models in the context of our class?}\spc{3}

\begin{tcolorbox}[colback=white, sharp corners]

This above excerpt means that even models that make correct predictions based upon training data cannot be fully accepted as true/correct because there may be a scenario or data point that exists for which the model may make an incorrect prediction. We can never know if a model makes correct predictions 100\% of the time, but we can be 100\% certain that when a model makes an incorrect prediction(s) that it would be wrong. 

\end{tcolorbox}

%--------------------------------------------
% Part (g)
%--------------------------------------------
\hardsubproblem{Further he says, \qu{you cannot prove a \emph{vague} theory wrong} (5:10 - 5:48). What does this mean in the context of mathematical models and metrics?} \spc{3}

\begin{tcolorbox}[colback=white, sharp corners]

If our mathematical models come with ill-defined metrics, then we cannot definitely make the claim that they are wrong. This idea makes me think about the logical implication that $p \implies q$ is true when $p$ is false but $q$ is true.

\end{tcolorbox}

%--------------------------------------------
% Part (h)
%--------------------------------------------
\hardsubproblem{He then continues with an example from psychology. Remember in the 1960's psychoanalysis was very popular. What is his remedy for being able to prove the vague psychology theory right (5:49 - 6:29)?} \spc{2}

\begin{tcolorbox}[colback=white, sharp corners]

Feynman's remedy for being able to prove the vague psychology theory right or wrong would be knowing ahead of time, using his psychology example, exactly how much ``love'' is not enough  and exactly how much ``love'' is over indulgent, we would then be able to test the theory and make definitive statements whether the model is wrong. 

\end{tcolorbox}

%--------------------------------------------
% Part (i)
%--------------------------------------------
\hardsubproblem{He then says \qu{then you can't claim to know anything about it} (6:40). Why can't you know anything about it?} \spc{3}

\begin{tcolorbox}[colback=white, sharp corners]

What Feynman is saying here is that ``you can't claim to know anything about'' the thing you do not clearly define the terms and conditions for. 

\end{tcolorbox}


Just to demonstrate that this modeling enterprise is all over science (not just Physics), I present to you the controversial theoretical political scientist John Mearsheimer. He's all over youtube and there's nothing special about this video that I will link here about \href{https://www.youtube.com/watch?v=D_Mx_e8t7nU&t=4673s}{Can China Rise Peacefully?} Feel free to watch the entire clip, but for the purposes of this class, we are only interested in the following segments referenced in the questions which has nothing to do with China, only his theory of \qu{power politics}.

%--------------------------------------------
% Part (j)
%--------------------------------------------
\hardsubproblem{Is Mearsheimer's model of great power politics / international relations (i.e., modern history) 9:35-17:22 simple or complicated? Explain.} \spc{3}

\begin{tcolorbox}[colback=white, sharp corners]

In my opinion, Mearsheimer's model of great power politics/international relations is a simple model based upon the five assumptions he presumes that the international political system adheres to. The phenomenon being modeled is complex because of the multiple actors that exist in said system and the wide spectrum of interactions that exist between the countries. The simplicity of his model lies in the fact that the number of assumptions he makes for his model is such a low number and he simplifies the model by making sweeping generalizations for every actor such as 1) no country knows the true intentions of another with any degree of certainty, and 2) survival is the only goal for each country. 

\end{tcolorbox}

%--------------------------------------------
% Part (k)
%--------------------------------------------
\hardsubproblem{Summarize his ideas about limitations of his theory from 39:18-40:00 using vocabulary from this class.} \spc{3}

\begin{tcolorbox}[colback=white, sharp corners]

Mearsheimer concedes that if his theory is considered to be ``one of the best theories'' about the phenomena he is modeling his theory would then be correct 75\% of the time and is incorrect 25\% of the time.  The limitations for the veracity of a theory is based upon two main things:

\begin{enumerate}
 \item[1)] the proxies $x_1, \ldots, x_n$, to the proximal drivers, $z_1, \ldots, z_t$, of the phenomena may be too much of an oversimplification, and 
 \item[2)] there may be proxies $x_{\cdot 1}, \ldots, x_{\cdot p}$ that are not even accounted for in the underlying assumptions for the model 
 \end{enumerate}

\end{tcolorbox}

\end{enumerate}

\end{document}

